% Nom : BELKASMI Abdeldjalil - ING4 %

\documentclass[12pt,a4paper]{article}
\usepackage[utf8]{inputenc}
\usepackage{enumitem}

\begin{document}

\section*{Exercice 1}

\indent Mes premiers pas avec \LaTeX{} sont un peu émouvants. Pour les premiers exemples, il faudra taper exactement ce qui est proposé, en faisant bien attention!

Les caractères spéciaux comme le pourcentage \%, le dollar \$, le souligné \_ ou l'esperluette \& nécessitent une attention particulière. Le projet est validé à 100\%.

\section*{Exercice 2}
Un texte contenant un mot en \textbf{gras}, un autre en \textit{italique}, et un terme technique en police \texttt{machine à écrire}.

Une phrase illustrant différentes tailles de caractères du texte normal, un mot en {\small petit}, un mot en {\large grand}, et un mot en {\Huge Énorme}.\\

Le premier paragraphe est aligné normalement sur les deux marges, ce qui correspond au comportement par défaut de \LaTeX{} pour tous les paragraphes du document.

\begin{flushleft}
Le deuxième paragraphe est aligné au fer à gauche, ce qui signifie que le bord droit est déchiqueté et non aligné.
\end{flushleft}

\begin{center}
Le troisième paragraphe est entièrement centré par rapport aux marges de la page.
\end{center}

\section*{Exercice 3}
Les domaines clés de l'intelligence artificielle sont :
\begin{enumerate}
    \item l'apprentissage automatique en raison:
    \begin{enumerate}
        \item de sa capacité de généralisation;
        \item de sa flexibilité face aux données massives;
        \item de ses performances en prédiction.
    \end{enumerate}
    \item le traitement du langage naturel parce que:
    \begin{enumerate}
        \item il permet aux machines de comprendre :
        \begin{itemize}[label=---]
            \item la sémantique textuelle;
            \item le contexte conversationnel;
            \item les nuances linguistiques.
        \end{itemize}
        \item il révolutionne les assistants intelligents.
    \end{enumerate}
\end{enumerate}

\end{document}
